\chapter{Conclusiones}
A partir de la información obtenida mediante el procedimiento experimental y el análisis de datos se llega a las siguientes conclusiones básicas:
\begin{itemize}
    \item Las sumatorias de torques obtenidas a través del procedimiento son consistentes con la hipótesis planteada y tienden a cero con las especificaciones dadas. Por lo tanto, se puede afirmar que los objetos se encuentran en un estado muy cercano al equilibrio rotacional.

    \item Las diferencias relativas respecto del equilibrio se encontraron entre 0,05\% y 3,47\%. Los casos <<e>> y <<h>> presentaron la mayor cercanía a la condición teórica, con diferencias relativas de 0,12\% y 0,05\% respectivamente. Estos resultados muestran que el modelo teórico de equilibrio rotacional fue compatible con las mediciones realizadas dentro del rango de masas y distancias estudiados.
    
    \item Las pequeñas diferencias respecto de la condición ideal $\sum \tau = 0$ pueden atribuirse a la resolución limitada de la balanza y de la cinta métrica, al rozamiento en el eje de suspensión, a la dificultad para mantener la viga completamente horizontal y a una posible distribución no uniforme de su masa. Por esta razón, los resultados no deben considerarse exactos, sino cercanos al modelo dentro de las limitaciones experimentales.
    
    \item En caso de repetir el procedimiento experimental se recomienda:
        \begin{enumerate}
            \item Mejorar, o re-diseñar el método usado para medir las masas, pues la balanza presentaba irregularidades al ser usada con masas de mayor orden.
            \item Derivado del anterior punto, tener un mejor control de las masas usadas, idealmente un mayor rango de masas para ajustar las masas hasta el punto deseado.
            \item Usar, en vez de una viga de madera, un objeto de masa homogénea y regular, como una varilla de metal.
        \end{enumerate}
\end{itemize}